\cotefile{79}
\scanpage{220}
qui \unsure{n'est autre|n'est-autre} que l'iso. \inserted{du type} (76), pour le système
dual du système envisagé \ldots

\deleted{\illegible{1} Il \unsure{reste donc qu'ainsi|reste donc que si l'on}}

\uline{Transport de structure}. Soit $g$ un autom. de
$(S,S^\vee,R)$, et soit $\varphi:S\longrightarrow S^\vee$ une jauge (il n'est
pas question d'un $Z$, pour l'instant). Alors \deleted{\illegible{1}}
$\varphi$ définit une structure de graphe $A_\varphi$, et on
a par transport de structure
\[
\left\{\begin{aligned}
g(\varphi)&\overset{\mathrm{déf}}=g_{S^\vee}\circ\varphi\circ g_S^{-1} &&\text{d'où :}\\
g_S(A_\varphi)&=A_{g(\varphi)}
\end{aligned}\right. \tag*{(87)}
\]

Quand on dispose d'un $Z$ ss-groupe comm. de
$\mathrm{Aut}(\Gamma=(S,S^\vee,\mathcal R))$, \inserted{$\varphi$ est une $g$-jauge, par ex. quand}
\inserted{et $g\in Z$, $\varphi_0\in\Pi_Z$, alors}
\deleted{du centralisateur $G$ dans $\mathrm{Aut}(\Gamma)$,}
\deleted{alors si $g\in G$ est une}
\deleted{alors l'ens. de $Z$-jauges $\varphi$ -- i.e. satisfaisant}
\deleted{$\lambda_{S^\vee}\varphi=\varphi\lambda_S^{-1}\quad\forall\lambda\in Z$, \unsure{déduction} \unsure{vide}, par une $G$-torsion via}
\[
 T_g\varphi_0=g_{S^\vee}\varphi_0=\varphi_0g_S^{-1} \tag*{(88)}
\]

\deleted{(NB que quelque soit $g$, si $Z$-jauges fixée, et $g\in\mathrm{Aut}(\Gamma)$,)}
\[
 \deleted{g_{S^\vee}\varphi=\varphi g_S^{-1}\quad\Longleftrightarrow\quad
 g\in\mathrm{Cent}_{\mathrm{Aut}(\Gamma)}(Z)\;(=G)} \tag*{\deleted{(89)}}
\]

Donc, on a alors, \deleted{pour $g\in G$}
\[
 \forall g\in\deleted{\unsure{G}}Z,\quad
 \varphi_0\in\underbrace{\Pi_Z}_{\substack{\text{ens. des}\\ Z\text{-jauges}}}
 \qquad \underbrace{g(\varphi)}_{\text{cf (87)}}=\deleted{\unsure{g}}T_{g^2}\varphi
 \tag*{(90)}
\]

Plus généralement (pour mémoire) si $\pi$ est un sous-$Z$-tors. de $\Pi_Z$,
\[
\begin{aligned}
&\text{Si }g\in G=\mathrm{Cent}_{\mathrm{Aut}(\Gamma)}(Z),\text{ alors conditions équivalentes}\\
&\left\{\begin{aligned}
\text{a)}\quad g_{S^\vee}\varphi&=\varphi g_S^{-1} &&\text{pour \uline{un} }\varphi\in\pi\quad(\text{disons, }\varphi_0)\\
\text{b)}\quad g_S\varphi&=\varphi g_{S^\vee}^{-1} &&\text{pour tt }\varphi\in\pi
\end{aligned}\right.\\
&\hspace{3em}(\text{NB relation qui s'écrit }\varphi(g_S)=(g_{S^\vee})^{-1}
\end{aligned}\tag*{(91)}
\]
% slanted
\margin{\uline{NB} Correspondent aux homothéties pour les structures $A_\varphi$ ($\varphi\in\pi$).
Les $g$ forment \deleted{sont \unsure{spéciales}} \unsure{un sous|un ensemble} \deleted{\unsure{stratification|stabilisateur}} (\unsure{i.e.|éqv.} sous \deleted{\unsure{Z}} s.g.) c'est pareil) sur $Z$}
et alors on \deleted{\unsure{a}} a
\[
 g(\varphi)=T_{g^2}\varphi\qquad\text{pour toute }\varphi\in\pi.
 \tag*{(92)}
\]
% archive sheet number: 219

